\documentclass[11pt,a4paper]{article}

% ── Packages ──────────────────────────────────────────────
\usepackage[utf8]{inputenc}
\usepackage[T1]{fontenc}
\usepackage{lmodern}
\usepackage{microtype}
\usepackage{amsmath,amssymb,amsthm}
\usepackage{mathtools}
\usepackage{geometry}
\usepackage{enumitem}
\usepackage{booktabs}
\usepackage{tabularx}
\usepackage{hyperref}
\usepackage{xcolor}
\usepackage[breakable,skins]{tcolorbox}
\usepackage{fancyhdr}
\usepackage{titlesec}
\usepackage{graphicx}
\usepackage{needspace}

\geometry{margin=2.5cm, headsep=14pt}

% ── Colors ────────────────────────────────────────────────
\definecolor{defblue}{RGB}{0,80,160}
\definecolor{keyred}{RGB}{180,30,30}
\definecolor{noteorange}{RGB}{200,120,0}
\definecolor{intugreen}{RGB}{30,130,60}
\definecolor{lightgray}{RGB}{245,245,245}
\definecolor{notabg}{RGB}{250,250,235}

% ── tcolorbox environments ────────────────────────────────
\newtcolorbox{keyidea}[1][]{%
  breakable, colback=red!5!white, colframe=keyred,
  fonttitle=\bfseries, title={Key Idea: #1}}

\newtcolorbox{intuition}[1][]{%
  breakable, colback=green!5!white, colframe=intugreen,
  fonttitle=\bfseries, title={Intuition: #1}}

\newtcolorbox{definition}[1][]{%
  breakable, colback=blue!5!white, colframe=defblue,
  fonttitle=\bfseries, title={Definition: #1}}

\newtcolorbox{warning}[1][]{%
  breakable, colback=orange!5!white, colframe=noteorange,
  fonttitle=\bfseries, title={Note: #1}}

\newtcolorbox{notationbox}[1][]{%
  breakable, colback=notabg, colframe=noteorange!70!black,
  fonttitle=\bfseries, title={Notation: #1}}

% ── Theorem-like environments ─────────────────────────────
\theoremstyle{plain}
\newtheorem{theorem}{Theorem}[section]
\newtheorem{lemma}[theorem]{Lemma}
\theoremstyle{remark}
\newtheorem*{remark}{Remark}

% ── Shortcuts (matching num_math notation) ────────────────
%    Basis functions : \psi_i           (Lectures 14--16)
%    Discrete space  : V_n              (Lecture 14)
%    Galerkin matrix : A_n              (Lecture 14)
%    Unknowns        : d_i              (Lecture 14)
%    Discrete soln   : u_n              (Lecture 14)
%    Test (discrete) : v_n              (Lecture 14)
%    Bilinear form   : a(\cdot,\cdot)   (all lectures)
%    Linear form     : \ell(\cdot) / F  (Lectures 14--15)
\newcommand{\R}{\mathbb{R}}
\newcommand{\N}{\mathbb{N}}
\newcommand{\dd}{\,\mathrm{d}}
\newcommand{\pd}[2]{\frac{\partial #1}{\partial #2}}
\newcommand{\norm}[1]{\left\| #1 \right\|}
\newcommand{\abs}[1]{\left| #1 \right|}
\newcommand{\inner}[2]{\left\langle #1,\, #2 \right\rangle}
\newcommand{\dx}{\,\mathrm{d}x}
\newcommand{\ds}{\,\mathrm{d}s}
\newcommand{\dA}{\,\mathrm{d}A}
\newcommand{\Vn}{V_n}
\newcommand{\un}{u_n}
\newcommand{\vn}{v_n}

% jump brackets (no extra package needed)
\newcommand{\jump}[1]{[\![#1]\!]}

% ── Header / Footer ───────────────────────────────────────
\pagestyle{fancy}
\fancyhf{}
\fancyhead[L]{\small FEM for Navier--Stokes\;---\;Study Guide}
\fancyhead[R]{\small SS\,2026}
\fancyfoot[C]{\thepage}
\renewcommand{\headrulewidth}{0.4pt}

% ── Title formatting ──────────────────────────────────────
% Each chapter (section) starts on a new page.
% Subsections and below keep at least 4 lines with their body.
\titleformat{\section}
  {\Large\bfseries\color{defblue}}{Chapter \thesection:}{0.5em}{}
\titlespacing*{\section}{0pt}{0pt}{6pt}

\titleformat{\subsection}
  {\large\bfseries\color{defblue!80!black}}{\thesubsection}{0.5em}{}
\titleformat{\subsubsection}
  {\normalsize\bfseries\color{defblue!60!black}}{\thesubsubsection}{0.5em}{}

% Prevent orphaned headings: keep subsection/subsubsection/paragraph
% glued to at least a few lines of the following content.
\titlespacing*{\subsection}{0pt}{12pt plus 2pt minus 2pt}{4pt}
\titlespacing*{\subsubsection}{0pt}{10pt plus 2pt minus 2pt}{3pt}

% Hook: before every \subsection, demand at least 6\baselineskip
\let\origsubsection\subsection
\renewcommand{\subsection}{\needspace{6\baselineskip}\origsubsection}
\let\origsubsubsection\subsubsection
\renewcommand{\subsubsection}{\needspace{5\baselineskip}\origsubsubsection}
\let\origparagraph\paragraph
\renewcommand{\paragraph}[1]{\needspace{4\baselineskip}\origparagraph{#1}}

\hypersetup{
  colorlinks=true, linkcolor=defblue, urlcolor=defblue!70!black
}

% ══════════════════════════════════════════════════════════
\begin{document}

\begin{center}
  {\LARGE\bfseries FEM for the Navier--Stokes Equations}\\[4pt]
  {\Large Personal Study Guide}\\[10pt]
  {\large Based on prerequisite lectures (Num.\ Math.\ II)
          and Chapter~2 slides}\\[6pt]
  {\normalsize SS\,2026 \quad$\cdot$\quad TU Berlin
               \quad$\cdot$\quad Dr.\ Peter M\"unch}\\[12pt]
  \rule{0.6\textwidth}{0.4pt}
\end{center}

\vspace{0.3cm}

\begin{notationbox}[Conventions used in this document]
This guide follows the notation of the \textbf{Num.\ Math.\ II} lectures
(Breiten \& May).  The Chapter~2 slides (M\"unch) use slightly different
symbols.  A translation table is provided in Section~\ref{sec:notation}.

\medskip
\renewcommand{\arraystretch}{1.2}
\begin{center}
\begin{tabular}{@{}lcc@{}}
  \toprule
  \textbf{Concept} & \textbf{Num.\ Math.\ II} & \textbf{Chapter~2 slides}\\
  \midrule
  Basis / shape functions & $\psi_i$       & $\varphi_j$     \\
  Discrete subspace       & $V_n$          & $V_h$            \\
  Discrete solution       & $u_n$          & $u_h$            \\
  Unknown coefficients    & $d_i$          & $U_j$            \\
  Galerkin / stiffness matrix & $A_n$      & $K$              \\
  Right-hand-side vector  & $\mathbf{f}$   & $F$              \\
  \bottomrule
\end{tabular}
\end{center}
\end{notationbox}

\vspace{0.2cm}
\tableofcontents
\newpage

% ══════════════════════════════════════════════════════════
\section{The Big Picture}
\label{sec:bigpicture}
% ══════════════════════════════════════════════════════════

The entire FEM pipeline answers one question:
\textbf{how do we turn a PDE (a continuous, infinite-dimensional problem)
into a linear system $A_n\,\mathbf{d} = \mathbf{f}$ that a computer can
solve?}

\begin{keyidea}[The FEM Pipeline]
\begin{center}
\small
$\displaystyle
  \underset{\text{\normalsize strong form}}{\boxed{\;-\Delta u = f\;}}
  \;\xrightarrow[\text{integrate}]{\text{multiply by }v}\;
  \underset{\text{\normalsize weak form}}{\boxed{\;a(u,v) = \ell(v)\;}}
$

\bigskip

$\displaystyle
  \xrightarrow[\text{dim.\ reduction}]{\text{Galerkin}}\;
  \underset{\text{\normalsize discrete form}}{\boxed{\;a(\un,\vn)=\ell(\vn)\;}}
  \;\xrightarrow[\text{$\un=\sum d_i\psi_i$}]{\text{basis expansion}}\;
  \underset{\text{\normalsize linear system}}{\boxed{\;A_n\,\mathbf{d}=\mathbf{f}\;}}
$
\end{center}
Every concept below is a piece of this chain.
\end{keyidea}

% ══════════════════════════════════════════════════════════
\newpage
\section{The Model Problem: Poisson Equation}
\label{sec:poisson}
% ══════════════════════════════════════════════════════════

\subsection{Strong form}

Let $\Omega \subset \R^d$ be a bounded domain with boundary
$\partial\Omega = \Gamma_D \cup \Gamma_N$, $\Gamma_D \cap \Gamma_N = \emptyset$.

\begin{definition}[Poisson equation --- strong form]
Find $u : \bar{\Omega} \to \R$ such that
\begin{alignat}{2}
  -\Delta u &= f        &\quad& \text{in } \Omega,
    \label{eq:poisson}\\
         u  &= g        &\quad& \text{on } \Gamma_D
    \quad\text{(Dirichlet b.c.)}, \label{eq:dirichlet}\\
  \pd{u}{n} &= h        &\quad& \text{on } \Gamma_N
    \quad\text{(Neumann b.c.)}. \label{eq:neumann}
\end{alignat}
Here $\Delta u = \sum_{i=1}^{d}\dfrac{\partial^2 u}{\partial x_i^2}$
is the Laplacian and
$\dfrac{\partial u}{\partial n} = \nabla u \cdot \mathbf{n}$
is the outward normal derivative.
\end{definition}

\begin{intuition}[What does the Poisson equation model?]
Think of $u$ as temperature.
The source term $f$ describes where heat is generated or absorbed.
The Laplacian $\Delta u$ measures how much $u$ at a point differs from
its average over nearby points.
The equation says: at steady state, the curvature of $u$ is balanced
by the source~$f$.

\medskip
\textbf{Dirichlet b.c.}\;($u = g$): we prescribe the \emph{value} on
part of the boundary (e.g.\ fixed temperature).\\[2pt]
\textbf{Neumann b.c.}\;($\partial u/\partial n = h$): we prescribe the
\emph{flux} on the rest (e.g.\ insulated wall $\Rightarrow h=0$).
\end{intuition}

% ══════════════════════════════════════════════════════════
\newpage
\section{From Strong to Weak Form}
\label{sec:weak}
% ══════════════════════════════════════════════════════════

\textit{Source: Num.\ Math.\ II, Lecture~12.}

\subsection{Why we need a weak form}

The strong form~\eqref{eq:poisson} requires
$u \in C^2(\Omega) \cap C(\bar{\Omega})$, i.e.\ $u$ must be twice
continuously differentiable.  This is too restrictive:
\begin{itemize}[nosep]
  \item Materials with discontinuous properties (e.g.\ a dielectric
        constant that jumps) produce solutions with ``kinks''
        (Lecture~15, num\_math).
  \item The weak form only needs \emph{first} derivatives, admitting
        rougher but physically meaningful solutions.
\end{itemize}

\subsection{Derivation (4 steps)}

\begin{definition}[Function spaces]
The \textbf{Sobolev space} $H^1(\Omega)$ contains all functions that are
square-integrable together with their first (weak) derivatives:
\[
  H^1(\Omega)
  = \bigl\{\, v : \Omega \to \R
    \;\big|\;
    v \in L^2(\Omega),\;
    \nabla v \in [L^2(\Omega)]^d
    \bigr\}.
\]
Important subspaces used in the lectures:
\begin{align*}
  C_0^2(\bar{\Omega})
    &= \bigl\{\, v \in C^2(\bar{\Omega})
       \;\big|\; v(0) = 0 = v(1) \bigr\}
    && \text{(Lecture~12: first test space)},\\[2pt]
  V_0
    &= \bigl\{\, v \in H^1(\Omega)
       \;\big|\; v = 0 \text{ on } \Gamma_D \bigr\}
    && \text{(general Dirichlet test space)},\\[2pt]
  V
    &= \bigl\{\, v \in H^1(\Omega)
       \;\big|\; v = 0 \text{ on } \Gamma_D \bigr\}
    && \text{(trial space, when $g=0$; otherwise lifted)}.
\end{align*}
\end{definition}

\medskip\noindent
Starting from $-\Delta u = f$ in $\Omega$:

\paragraph{Step 1: Multiply by a test function and integrate.}
Let $v \in V_0$ (so $v = 0$ on $\Gamma_D$).
Multiply both sides by~$v$ and integrate over~$\Omega$:
\[
  \int_\Omega (-\Delta u)\, v \dx
  = \int_\Omega f\, v \dx.
\]

\paragraph{Step 2: Integration by parts (Green's first identity).}

\begin{keyidea}[The central trick]
Integration by parts shifts one derivative from $u$ onto $v$.
Recall the multivariate product rule (Lecture~12, p.\,4):
$\operatorname{div}(v\,\nabla u) = \nabla u \cdot \nabla v + v\,\Delta u$.
Integrating and applying Gauss's theorem:
\[
  \int_\Omega (-\Delta u)\, v \dx
  \;=\;
  \underbrace{\int_\Omega \nabla u \cdot \nabla v \dx}_{%
    \text{only first derivatives!}}
  \;-\;
  \underbrace{\int_{\partial\Omega}
    v \cdot \pd{u}{n} \dA}_{\text{boundary integral}}.
\]
This is why the weak form needs only first derivatives of both $u$
and $v$.
\end{keyidea}

\paragraph{Step 3: Apply boundary conditions.}
Split the boundary integral over $\Gamma_D$ and $\Gamma_N$:
\begin{itemize}[nosep]
  \item On $\Gamma_D$: $v = 0$ by choice of test space
        $\;\Rightarrow\;$ integral vanishes.
  \item On $\Gamma_N$: $\partial u/\partial n = h$ (Neumann data)
        $\;\Rightarrow\;$ integral becomes
        $\displaystyle\int_{\Gamma_N} h\,v\ds$.
\end{itemize}
(Lecture~12, p.\,4 shows this step-by-step for the 2D case with
$\Gamma = \Gamma_1 \cup \Gamma_2$.)

\paragraph{Step 4: Write the weak form.}

\begin{definition}[Weak / variational formulation\; (VP)]
Find $u \in V$ such that
\begin{equation}
  \boxed{\;a(u,v) = \ell(v) \qquad \forall\, v \in V_0\;}
  \label{eq:weak}
\end{equation}
where the \textbf{bilinear form} and \textbf{linear functional} are:
\begin{align}
  a(u,v)  &= \int_\Omega \nabla u \cdot \nabla v \dx,
             \label{eq:bilinear}\\[4pt]
  \ell(v) &= \int_\Omega f\, v \dx
             + \int_{\Gamma_N} h\, v \ds.
             \label{eq:linear}
\end{align}
\end{definition}

\begin{intuition}[Why does weak\,$=$\,strong?]
Lecture~12 proves this using the \textbf{Fundamental Lemma of Calculus of
Variations}: if\, $\int_\Omega g(x)\,h(x)\dx = 0$\, for all smooth $h$
with compact support, then $g \equiv 0$.

Applied here: define $g(x) = -u''(x) - f(x)$.  The weak form says
$\int g\,v\dx = 0$ for all $v \in C_0^\infty$.  By the lemma,
$g \equiv 0$, i.e.\ $-u'' = f$ pointwise---recovering the strong form.
So the two formulations are \textbf{equivalent} whenever $u$ is smooth
enough.
\end{intuition}

\begin{warning}[Neumann b.c.\ are ``natural'']
Neumann conditions appear \emph{automatically} in the weak form via the
boundary integral---you do \emph{not} enforce them explicitly.
Dirichlet conditions, by contrast, are \emph{essential}: they are
enforced by restricting the function space
(Lecture~17, num\_math).
\end{warning}

% ══════════════════════════════════════════════════════════
\newpage
\section{The Galerkin Method}
\label{sec:galerkin}
% ══════════════════════════════════════════════════════════

\textit{Source: Num.\ Math.\ II, Lecture~14.}

\subsection{Three-step recipe}

\paragraph{Step 1 (known): Find the variational / weak form.}
Find $u \in V$ such that $a(u,v) = \ell(v)$ for all $v \in V$.
The space $V$ is infinite-dimensional---we cannot work with it directly
on a computer.

\paragraph{Step 2 (new): Galerkin dimension reduction.}
For each $n \in \N$, let $\Vn \subset V$ be a subspace with
$\dim(\Vn) = n$.  Pick a basis $\{\psi_1, \ldots, \psi_n\}$ of $\Vn$.

\begin{definition}[Galerkin equation]
Find $\un \in \Vn$ such that
\[
  a(\un, \vn) = \ell(\vn)
  \qquad \forall\, \vn \in V_{0,n} \subset V_0.
\]
\end{definition}

\paragraph{Step 3 (new): Derive a linear system of equations.}
Since $\Vn$ is finite-dimensional, every element can be written as
\[
  \un(x) = \sum_{i=1}^{n} d_i\,\psi_i(x),
\]
where $d_1, \ldots, d_n \in \R$ are the unknown coefficients.

\begin{keyidea}[Testing with basis functions $\psi_j$ is enough]
Instead of testing with arbitrary $\vn \in \Vn$, it suffices to test
with each basis function $\psi_j$, $j = 1, \ldots, n$.

\medskip\noindent
\textit{Why?}\; Any $\vn = \sum_{j} \beta_j\,\psi_j$.
By linearity of $a$ and $\ell$:
\[
  a(\un,\vn) = \ell(\vn)
  \;\;\Longrightarrow\;\;
  \sum_j \beta_j\bigl[a(\un,\psi_j) - \ell(\psi_j)\bigr] = 0.
\]
Since the $\beta_j$ are arbitrary, each bracket must vanish individually.
\end{keyidea}

Substituting $\un = \sum_{i} d_i\,\psi_i$ and testing with $\psi_j$:
\[
  a\!\Bigl(\sum_{i=1}^n d_i\,\psi_i,\;\psi_j\Bigr)
  = \ell(\psi_j)
  \qquad\xRightarrow{\text{bilinearity}}\qquad
  \sum_{i=1}^n d_i\, a(\psi_i,\psi_j) = \ell(\psi_j).
\]
This is a linear system of $n$ equations for $n$ unknowns:

\begin{equation}
  \boxed{\;A_n\,\mathbf{d} = \mathbf{f}\;}
  \label{eq:system}
\end{equation}
with
\begin{align}
  (A_n)_{ji}
    &= a(\psi_i,\, \psi_j)
     = \int_\Omega \nabla\psi_i \cdot \nabla\psi_j \dx,
    \label{eq:stiffness}\\[4pt]
  f_j
    &= \ell(\psi_j)
     = \int_\Omega f\,\psi_j \dx
       + \int_{\Gamma_N} h\,\psi_j \ds.
    \label{eq:load}
\end{align}

\medskip
$A_n \in \R^{n \times n}$ is the \textbf{Galerkin matrix}
(also called \textbf{stiffness matrix}),
$\mathbf{f} \in \R^n$ is the \textbf{right-hand-side (load) vector},
and $\mathbf{d} = (d_1,\ldots,d_n)^T$ are the unknown
\textbf{degrees of freedom} (DOFs).

\subsection{Galerkin orthogonality and best approximation}

\textit{Source: Num.\ Math.\ II, Lecture~14, p.\,2.}

\medskip\noindent
Since the true solution $u$ of~(VP) satisfies $a(u,v) = \ell(v)$ for all
$v \in V$, in particular for all $\vn \in \Vn \subset V$:
\[
  a(u,\vn) = \ell(\vn) = a(\un,\vn)
  \qquad\forall\,\vn \in \Vn.
\]
Subtracting:

\begin{definition}[Galerkin orthogonality]
\begin{equation}
  \boxed{\;a(u - \un,\;\vn) = 0
         \qquad \forall\, \vn \in \Vn\;}
  \label{eq:galerkin-orth}
\end{equation}
The error $u - \un$ is ``$a$-orthogonal'' to the entire discrete
subspace~$\Vn$.
\end{definition}

\begin{intuition}[Best approximation]
If $a$ is symmetric and positive definite, Galerkin orthogonality means
$\un$ is the \textbf{closest} element in $\Vn$ to the true solution $u$,
measured in the \textbf{energy norm}
$\norm{v}_a = \sqrt{a(v,v)}$:
\[
  \norm{u - \un}_a
  = \inf_{\vn \in \Vn} \norm{u - \vn}_a.
\]
The FEM does not give just \emph{some} approximation---it gives the
\textbf{best possible} one from the chosen subspace.
\end{intuition}

\subsection{Properties of the bilinear form}

\textit{Source: Num.\ Math.\ II, Lecture~14, p.\,2.}

\begin{definition}[Bilinear form properties]
Let $V$ be an $\R$-vector space and $a : V \times V \to \R$ a bilinear
form.  We call~$a$:
\begin{itemize}[nosep]
  \item \textbf{symmetric}, if $a(u,v) = a(v,u)$ for all $u,v \in V$;
  \item \textbf{positive definite}, if $a(u,u) \geq 0$ for all $u$
        and $a(u,u) = 0$ if and only if $u = 0$;
  \item \textbf{negative definite}, if $-a$ is positive definite.
\end{itemize}
If $a$ is positive definite, the variational problem has a
\textbf{unique} solution.
\end{definition}

\begin{warning}[Positive definiteness depends on the space!]
Lecture~14 (p.\,3) gives a crucial example:
$a(u,v) = \int_\Omega \nabla u \cdot \nabla v\dx$.
\begin{enumerate}[nosep]
  \item On $V = \{v \in H^1(\Omega) \mid \partial v/\partial n = 0
        \text{ on } \partial\Omega\}$:
        $a$ is \textbf{not} positive definite
        (constant functions have $\nabla u = 0$ but $u \neq 0$).
  \item On $V_0 = \{v \in H^1(\Omega) \mid v = 0
        \text{ on } \partial\Omega\}$:
        $a$ \textbf{is} positive definite
        (Dirichlet b.c.\ kill the constants).
\end{enumerate}
The choice of boundary conditions determines whether the problem is
well-posed!
\end{warning}

% ══════════════════════════════════════════════════════════
\newpage
\section{Finite Elements: The Choice of $V_n$}
\label{sec:elements}
% ══════════════════════════════════════════════════════════

\textit{Source: Num.\ Math.\ II, Lectures~15--16.}

\medskip\noindent
The Galerkin method works for any subspace $\Vn$.
The key question is: how to choose $\Vn$ (and hence $\psi_1,\ldots,\psi_n$)
so that (a)~the approximation is accurate, (b)~$A_n$ is sparse and
easy to assemble, and (c)~$\Vn$ behaves well as $n \to \infty$?

The answer from Lecture~15: \textbf{hat functions} (piecewise linear
polynomials on a mesh).

\subsection{Admissible decomposition (mesh)}

\begin{definition}[Admissible decomposition]
Decompose $\bar{\Omega} = \bigcup_{i=1}^{n_e} \bar{\Omega}_i$ into
elements.  The decomposition is \textbf{admissible} if for each pair
$i \neq j$ exactly one of the following holds:
\begin{enumerate}[nosep]
  \item $\bar{\Omega}_i \cap \bar{\Omega}_j = \emptyset$ \quad(disjoint),
  \item $\bar{\Omega}_i \cap \bar{\Omega}_j$ is a complete shared
        face/edge/vertex.
\end{enumerate}
\begin{itemize}[nosep]
  \item In 1D ($d=1$): non-overlapping intervals
        $\Omega_\ell = (x_{\ell-1}, x_\ell)$,
        $\ell = 1,\ldots,n_e = N+1$.
  \item In 2D ($d=2$): triangles or convex quadrilaterals with
        compatible edges (Lecture~15, p.\,3 illustrates admissible
        vs.\ non-admissible decompositions).
\end{itemize}
\end{definition}

\subsection{Shape functions (basis functions $\psi_i$)}

\subsubsection{1D: hat functions (Lectures~15--16)}

For a 1D mesh $a = x_0 < x_1 < \cdots < x_N < x_{N+1} = b$, the
\textbf{linear Lagrange basis} (``hat functions'' / ``pyramid functions'')
are:
\[
  \psi_i(x) = \begin{cases}
    \dfrac{x - x_{i-1}}{x_i - x_{i-1}}
      & x \in [x_{i-1},\, x_i],\\[8pt]
    \dfrac{x_{i+1} - x}{x_{i+1} - x_i}
      & x \in [x_i,\, x_{i+1}],\\[8pt]
    0 & \text{otherwise},
  \end{cases}
  \qquad i = 1,\ldots,N.
\]

Properties (Lecture~16, p.\,1):
\begin{itemize}[nosep]
  \item $\psi_i(x_j) = \delta_{ij}$
        (equals~1 at node~$i$, zero at all other nodes).
  \item $\psi_i$ is piecewise linear and continuous;
        $\psi_i \in H^1(\Omega)$ but $\psi_i \notin C^1(\bar\Omega)$.
  \item Each $\psi_i$ has \textbf{local support}: nonzero only on the
        two neighboring elements.
  \item $\Rightarrow\;(A_n)_{ij} = a(\psi_i,\psi_j) = 0$
        whenever $|i - j| > 1$
        $\;\Rightarrow\; A_n$ is \textbf{tridiagonal}.
\end{itemize}

\subsubsection{Explicit stiffness matrix
              (1D, equidistant mesh, homogeneous Dirichlet)}

\textit{Source: Num.\ Math.\ II, Lecture~16, p.\,1.}

\medskip\noindent
For $\Omega = (0,1)$, uniform mesh width $h = 1/(N+1)$, and $u(0) = u(1) = 0$:
\[
  (A_n)_{ij}
  = \int_0^1 \psi_i'(x)\,\psi_j'(x)\dx
  = \frac{1}{h}
  \begin{cases}
    \phantom{-}2  & \text{if } i = j,\\
    -1            & \text{if } |i-j| = 1,\\
    \phantom{-}0  & \text{otherwise}.
  \end{cases}
\]
\[
  A_n = \frac{1}{h}\begin{pmatrix}
    2 & -1 & & \\
    -1 & 2 & -1 & \\
    & \ddots & \ddots & \ddots \\
    & & -1 & 2
  \end{pmatrix} \in \R^{N \times N}.
\]
The right-hand side is typically computed by numerical quadrature:
$f_j = \ell(\psi_j) = \int_0^1 f(x)\,\psi_j(x)\dx
\approx \sum_q w_q\, f(x_q)\,\psi_j(x_q)$.

\begin{warning}[Same as finite differences!]
The FEM matrix $A_n = \frac{1}{h}\operatorname{tridiag}(-1,2,-1)$ is
identical to the finite-difference matrix from Lecture~01.
For linear elements in 1D the two methods coincide.
This is \emph{not} true in 2D or for higher-order elements.
\end{warning}

\subsubsection{2D element catalog (Lecture~16, pp.\,2--3)}

On each element $\bar{\Omega}_i$, the local basis function is a polynomial.
Common FEM choices:

\renewcommand{\arraystretch}{1.4}
\begin{center}
\begin{tabular}{@{}llcc@{}}
  \toprule
  \textbf{Element type} & \textbf{Cell shape}
    & \textbf{Nodes} & \textbf{DOFs/elem.}\\
  \midrule
  Linear ($P_1$)      & triangle       & 3  & 3  \\
  Quadratic ($P_2$)   & triangle       & 6  & 6  \\
  Cubic ($P_3$)       & triangle       & 10 & 10 \\
  Bilinear ($Q_1$)    & quadrilateral  & 4  & 4  \\
  Biquadratic ($Q_2$) & quadrilateral  & 9  & 9  \\
  \bottomrule
\end{tabular}
\end{center}

\begin{itemize}[nosep]
  \item Triangle with degree~$p$:
        $\binom{p+2}{2}$ nodes
        (polynomial $\psi_i|_{\bar\Omega_i}
        = a_0 + a_1 x_1 + a_2 x_2 + \cdots$).
  \item Quadrilateral with degree~$p$:
        $(p+1)^2$ nodes (tensor-product polynomial).
  \item The \textbf{Lagrange polynomial} formula
        $\ell_i(x) = \prod_{k \neq i}
        \frac{x - c_k}{c_i - c_k}$
        guarantees $\ell_i(c_j) = \delta_{ij}$
        (Lecture~16, p.\,2).
\end{itemize}

\begin{warning}[Dimension of $V_n$ $\neq$ DOFs/element $\times$
number of elements]
The dimension of the resulting subspace $V_n$ depends on the
\emph{global} node count and the boundary conditions, not simply on
the per-element DOF count (Lecture~16, p.\,3).
\end{warning}

% ══════════════════════════════════════════════════════════
\newpage
\section{Assembly: Building $A_n$ and $\mathbf{f}$}
\label{sec:assembly}
% ══════════════════════════════════════════════════════════

\subsection{Element-wise assembly}

We never compute $(A_n)_{ij}$ by integrating over all of~$\Omega$ at once.
Since shape functions have local support, each integral reduces to a
sum over element contributions:

\begin{keyidea}[Assembly via element matrices]
\[
  A_n = \sum_c R_c^T\, A_c\, R_c,
  \qquad
  \mathbf{f} = \sum_c R_c^T\, \mathbf{f}_c,
\]
where:
\begin{itemize}[nosep]
  \item $c$ loops over all cells (elements) in the mesh,
  \item $A_c$ is the \textbf{element stiffness matrix}
        (small, e.g.\ $3\times 3$ for a $P_1$ triangle),
  \item $\mathbf{f}_c$ is the \textbf{element load vector},
  \item $R_c$ is the \textbf{restriction matrix} mapping local DOF
        indices to global ones.
\end{itemize}
\end{keyidea}

\subsection{Integration by numerical quadrature}

The integrals
$(A_c)_{ij} = \int_{\Omega_c} \nabla\psi_j \cdot \nabla\psi_i\dx$
are evaluated numerically on a reference element~$\hat{\Omega}$:
\[
  \int_{\Omega_c} g(x)\dx
  \;\approx\;
  \sum_{q=1}^{n_q} w_q\, g(\mathbf{x}_q)\,|\det J|,
\]
where $\{\mathbf{x}_q, w_q\}$ are the quadrature points and weights,
and $J$ is the Jacobian of the mapping
$\hat{\Omega} \to \Omega_c$.

\subsection{Boundary conditions in practice}

\textbf{Neumann:} Already incorporated in $\ell(v)$---nothing extra
to do.

\medskip
\textbf{Dirichlet ($u = g$ on $\Gamma_D$):}
For each DOF $i$ on $\Gamma_D$:
\begin{enumerate}[nosep]
  \item Set row~$i$ of~$A_n$ to zero except $(A_n)_{ii} = 1$.
  \item Modify the other rows:
        $f_j \leftarrow f_j - (A_n)_{ji}\,g_i$ for each $j \neq i$.
  \item Set $f_i = g_i$.
\end{enumerate}
This enforces $d_i = g_i$ while preserving the solution at
interior nodes.

% ══════════════════════════════════════════════════════════
\newpage
\section{Error Estimates}
\label{sec:error}
% ══════════════════════════════════════════════════════════

\subsection{Interpolation error}

\begin{definition}[Polynomial interpolation error]
Given $f \in C^{n+1}[a,b]$ and its interpolating polynomial $p_n$ of
degree~$n$ at $n+1$ distinct points:
\[
  e_n(x)
  = f(x) - p_n(x)
  = \frac{f^{(n+1)}(\xi)}{(n+1)!}\;\omega_{n+1}(x),
  \quad
  \omega_{n+1}(x)
  = \prod_{k=0}^{n}(x - x_k).
\]
\end{definition}

For FEM shape functions of polynomial degree~$p$ on elements of
size~$h$:
\[
  \norm{\nabla e}_{L^2} \leq \frac{C(p)}{p!}\, h^p.
\]

\subsection{A priori error estimate}

Combining the interpolation error bound with Galerkin
orthogonality~\eqref{eq:galerkin-orth}:

\begin{keyidea}[A priori convergence rate]
For polynomial degree~$p$ and mesh size~$h$:
\begin{equation}
  \boxed{\;\norm{\nabla(u - \un)}_{L^2}
         \leq C\, h^p\, \norm{\nabla^{p+1} u}_{L^2}\;}
  \label{eq:apriori}
\end{equation}
\begin{itemize}[nosep]
  \item $p=1$ (linear elements): error $\sim \mathcal{O}(h)$.
        Halve the mesh $\Rightarrow$ halve the error.
  \item $p=2$ (quadratic elements): error $\sim \mathcal{O}(h^2)$.
        Halve the mesh $\Rightarrow$ quarter the error.
\end{itemize}
\end{keyidea}

\begin{warning}[Regularity matters]
The estimate requires $u \in H^{p+1}$ (the solution must be smooth
enough).  On domains with re-entrant corners (e.g.\ L-shaped), the
solution may have singularities that limit the actual convergence rate
regardless of~$p$---motivating adaptive mesh refinement.
\end{warning}

\subsection{A posteriori error estimate}

Given a \emph{computed} solution $\un$, estimate the error
\emph{without} knowing~$u$:
\begin{equation}
  \norm{e_n}_A
  \leq C \sqrt{\sum_T
    \Bigl(
      h_T^2\,\norm{R_T}_T^2
      + h_T\,\norm{J_E}_{\partial T}^2
    \Bigr)},
  \label{eq:aposteriori}
\end{equation}
where:
\begin{itemize}[nosep]
  \item $R_K = f + \Delta \un$ is the \textbf{cell residual}
        (how well $\un$ satisfies the PDE inside each element),
  \item $J_E = -\tfrac{1}{2}\jump{\partial_n \un}$ is the
        \textbf{edge residual} (jump in normal derivative across
        element boundaries).
\end{itemize}

This drives \textbf{adaptive mesh refinement (AMR)}:
\[
  \text{solve}
  \;\to\;
  \text{mark cells with large error indicators}
  \;\to\;
  \text{refine}
  \;\to\;
  \text{repeat}.
\]

% ══════════════════════════════════════════════════════════
\newpage
\section{Theory: Existence and Uniqueness}
\label{sec:laxmilgram}
% ══════════════════════════════════════════════════════════

\textit{Source: Num.\ Math.\ II, Lecture~17.}

\begin{definition}[Continuity and coercivity]
Let $\Omega \subset \R^d$ be a bounded polygonal domain and set
$V = H_0^1(\Omega)$ with norm
$\norm{u}_V = \norm{u}_{H^1(\Omega)}$.
Assume $a: V \times V \to \R$ and
$\ell: V \to \R$ satisfy:
\begin{enumerate}[label=\alph*),nosep]
  \item \textbf{Continuity of $a$:}\;
        $\exists\;\alpha > 0$ such that
        $|a(v_1,v_2)| \leq \alpha\,\norm{v_1}_V\,\norm{v_2}_V$
        $\;\forall\, v_1,v_2 \in V$.
  \item \textbf{Coercivity ($V$-ellipticity) of $a$:}\;
        $\exists\;\beta > 0$ such that
        $a(v,v) \geq \beta\,\norm{v}_V^2$
        $\;\forall\, v \in V$.
  \item \textbf{Continuity of $\ell$:}\;
        $\exists\;g > 0$ such that
        $|\ell(v)| \leq g\,\norm{v}_V$
        $\;\forall\, v \in V$.
\end{enumerate}
\end{definition}

\begin{theorem}[Lax--Milgram]
Under assumptions (a)--(c), the variational
problem~\eqref{eq:weak} has a \textbf{unique} solution $u \in V$ with
the stability estimate
\[
  \norm{u}_V \leq \frac{g}{\beta}.
\]
\end{theorem}

\begin{intuition}[Lax--Milgram $\approx$ invertibility in disguise]
Compare with linear algebra (Lecture~17, p.\,2):
\begin{center}
\renewcommand{\arraystretch}{1.2}
\begin{tabular}{@{}lcc@{}}
  \toprule
  & \textbf{Linear algebra} & \textbf{Variational problem}\\
  \midrule
  Problem    & $Ax = b$            & $a(u,v) = \ell(v)$\\
  Bilinear   & $a(x,y) = y^T A x$ & $a(u,v)$\\
  Continuity & $|y^T A x| \leq \norm{A}\norm{x}\norm{y}$
             & $|a(u,v)| \leq \alpha\norm{u}\norm{v}$\\
  Coercivity & $x^T A x \geq \beta \norm{x}^2$ (pos.\ def.)
             & $a(v,v) \geq \beta\norm{v}^2$\\
  Solution   & $x = A^{-1}b$, $\norm{x} \leq \norm{A^{-1}}\norm{b}$
             & $\norm{u}_V \leq g/\beta$\\
  \bottomrule
\end{tabular}
\end{center}
Coercivity is the infinite-dimensional version of positive definiteness.
\end{intuition}

Note: $a$ is \emph{not required to be symmetric} for Lax--Milgram
(Lecture~17, p.\,2).

\subsection{Variational $\Leftrightarrow$ Minimization
            (when $a$ is symmetric)}

\textit{Source: Num.\ Math.\ II, Lecture~15, p.\,1.}

\begin{theorem}[Equivalence of VP and MP]
Let $a$ be symmetric and positive definite.  Define
$\mathcal{J}: V \to \R$ by
\[
  \mathcal{J}(v) = \tfrac{1}{2}\,a(v,v) - \ell(v).
\]
Then the variational problem (VP) and the minimization problem
\[
  \text{(MP)}\quad
  \text{Find } u_{\min} \in V \text{ such that }
  \mathcal{J}(u_{\min}) \leq \mathcal{J}(v)
  \;\;\forall\, v \in V
\]
have the same (unique) solution.
\end{theorem}

\noindent
\textit{Proof sketch} (Lecture~15):
(VP)$\Rightarrow$(MP): Define $w = u_{\min} + \epsilon\,v$ and compute
$\tilde{\mathcal{J}}(\epsilon) = \mathcal{J}(u_{\min}+\epsilon v)$.
Minimum at $\epsilon = 0$ implies $\tilde{\mathcal{J}}'(0) = 0$,
which yields $a(u_{\min},v) = \ell(v)$.
The reverse direction follows similarly.

% ══════════════════════════════════════════════════════════
\newpage
\section{Extension to Higher Dimensions}
\label{sec:higherdim}
% ══════════════════════════════════════════════════════════

The variational framework is \textbf{dimension-independent}.
Moving from 1D to 2D/3D requires changing:
\begin{enumerate}[nosep]
  \item \textbf{Mesh:} triangles / quadrilaterals (2D),
        tetrahedra / hexahedra (3D).
  \item \textbf{Shape functions $\psi_i$:} e.g.\ linear on triangles
        (3~nodes), bilinear on quads (4~nodes).
  \item \textbf{Quadrature rules:} for numerical integration on the
        reference element.
  \item \textbf{Reference-to-physical mapping:}
        Jacobian~$J$ for coordinate transformation.
\end{enumerate}

Key consequence: in 2D with unstructured meshes, $A_n$ is still
\textbf{sparse} (each $\psi_i$ overlaps only with neighbors), but the
sparsity pattern is less regular than the clean banded structure in 1D.

% ══════════════════════════════════════════════════════════
\newpage
\section{Notation Reference}
\label{sec:notation}
% ══════════════════════════════════════════════════════════

\renewcommand{\arraystretch}{1.3}
\begin{center}
\begin{tabularx}{\textwidth}{@{}l l X@{}}
  \toprule
  \textbf{Term}
    & \textbf{Symbol}
    & \textbf{Meaning} \\
  \midrule
  Trial space
    & $V$
    & Space where we seek $u$
      (infinite-dim.) \\
  Test space
    & $V_0$
    & Functions $v$ we multiply by;
      $v=0$ on $\Gamma_D$ \\
  Discrete subspace
    & $V_n$
    & Finite-dim.\ subspace of $V$;
      $\dim(V_n) = n$ \\
  Basis functions
    & $\psi_1,\ldots,\psi_n$
    & Piecewise polynomials; basis of $V_n$ \\
  Unknown coefficients
    & $d_1,\ldots,d_n$
    & DOFs: $\un = \sum d_i\,\psi_i$ \\
  Discrete solution
    & $u_n$
    & Approximation in $V_n$ \\
  Bilinear form
    & $a(\cdot,\cdot)$
    & ``Stiffness'' part:
      $\int \nabla u \cdot \nabla v\dx$ for Poisson \\
  Linear form
    & $\ell(\cdot)$
    & ``Load'' part:
      $\int f\,v\dx + \int_{\Gamma_N} h\,v\ds$ \\
  Galerkin matrix
    & $A_n$
    & $(A_n)_{ji} = a(\psi_i,\,\psi_j)$;
      stiffness matrix \\
  Mass matrix
    & $M_n$
    & $(M_n)_{ij} = \int \psi_i\,\psi_j\dx$;
      appears in time-dep.\ problems \\
  RHS vector
    & $\mathbf{f}$
    & $f_j = \ell(\psi_j)$;
      load vector \\
  Galerkin orthogonality
    &
    & $a(u-\un,\,\vn) = 0\;\;\forall\,\vn \in V_n$ \\
  Lax--Milgram
    &
    & Continuity + coercivity
      $\Rightarrow$ unique solution \\
  A priori estimate
    &
    & Error bound in terms of $h$, $p$,
      and regularity of $u$ \\
  A posteriori estimate
    &
    & Error bound computable from $\un$
      alone \\
  AMR
    &
    & Adaptive mesh refinement:
      solve $\to$ estimate $\to$ refine $\to$ repeat \\
  \bottomrule
\end{tabularx}
\end{center}

% ══════════════════════════════════════════════════════════
\newpage
\section{Study Roadmap}
\label{sec:roadmap}
% ══════════════════════════════════════════════════════════

Recommended reading order for the prerequisite lectures:

\begin{enumerate}
  \item \textbf{Lecture~12} (pp.\,1--4):
        Weak form derivation---understand every line, especially the
        integration-by-parts step on p.\,4.
  \item \textbf{Lecture~14} (pp.\,1--2):
        Galerkin 3-step recipe ($V \to V_n \to A_n\mathbf{d}=\mathbf{f}$).
        Galerkin orthogonality and uniqueness.
  \item \textbf{Lecture~14} (p.\,3):
        Positive definiteness depends on the space---work through the
        two examples.
  \item \textbf{Lecture~15} (last 2~pages):
        Admissible decompositions and hat functions---the bridge from
        Galerkin to FEM.
  \item \textbf{Lecture~16} (p.\,1):
        Explicit 1D stiffness matrix---compute $A_n$ by hand for $N=3$.
  \item \textbf{Lecture~16} (pp.\,2--3):
        2D element catalog: $P_1$, $P_2$, $Q_1$, $Q_2$, Lagrange
        polynomials.
  \item \textbf{Lecture~17} (p.\,1):
        Neumann BCs in weak form.
  \item \textbf{Lecture~17} (pp.\,2--3):
        Lax--Milgram theorem and the analogy with linear algebra.
\end{enumerate}

\medskip
\textbf{Self-test:} Starting from $-\Delta u = f$, can you explain why
``multiply by~$v$ and integrate by parts'' ultimately produces the
matrix equation $A_n\,\mathbf{d} = \mathbf{f}$?  If yes, you are ready
for the lecture.

% ══════════════════════════════════════════════════════════
%
%  ── SPACE FOR FUTURE CHAPTERS ───────────────────────────
%  Add new \section{} blocks below as the course progresses.
%  Suggested upcoming sections:
%
%  \section{Software Development with Linux and C/C++}
%  \section{Elliptic \& Parabolic Problems}
%  \section{Hyperbolic Problems}
%  \section{Linear Solvers \& Multigrid}
%  \section{Stokes Problem}
%  \section{Navier--Stokes Problem}
%  \section{Adaptive Mesh Refinement}
%  \section{High-Performance Computing}
%
% ══════════════════════════════════════════════════════════

\end{document}
